\section{The monochromatic field with a Gaussian turn-on}
We study the motion of a charged particle in a plane wave field with a Gaussian turn-on.\\
We define the envelope as
\begin{align}
    f(\chi)= \begin{cases}e^{-\frac{\chi^{2}}{(2 \pi \sigma)^{2}}}, & \chi<\chi_{0}  \label{I.1.3.1}\\ 1, & \chi \geq \chi_{0}\end{cases}
\end{align}
and the vector potential in the Lorentz-Coulomb gauge as
\begin{align}
    \mathbf{A}(\phi)=A_{0} f(k \phi)\left(\mathbf{e}_{x} \zeta_{x} \cos \left(k \phi+\Phi_{0}\right)+\mathbf{e}_{y} \zeta_{y} \sin \left(k \phi+\Phi_{0}\right)\right) . \label{I.1.3.2}
\end{align}
where $k$ is the wave number $k=\frac{\omega}{c}, \sigma$ gives the width of the Gaussian envelope, $\Phi_{0}$ is the initial phase and $\zeta_{x}, \zeta_{y}$ are the polarization amplitudes, obeying the condition $\zeta_{x}^{2}+\zeta_{y}^{2}=1$.
It is convenient to define the unity vector of the field propagation direction as $\mathbf{n}=(0,0,1)$ and the four-vector
\begin{align}
    n=(1, \mathbf{n}) \label{I.1.3.3}
\end{align}
which in the case of the propagation along the $O z$ axis becomes $n=(1,0,0,1)$. Also the notation for the dimensionless field amplitude is
\begin{align}
    \xi_{0}=\frac{e_{0} A_{0}}{m c} . \label{I.1.3.4}
\end{align}
